\newif\ifuseseminar
\useseminartrue
\input{../../latex_common/header.tex.frag}

\title{Avaliação Diagnóstica}

% ------------------------------------------------------------------
% Configurações de página e estilo
\geometry{margin=2.5cm}
\hypersetup{hidelinks}
\pagestyle{fancy}
\fancyhf{}
\lhead{\small Universidade / Curso}
\cfoot{\small Página \thepage~de \pageref{LastPage}}

\setlength{\parindent}{0pt}
\setlength{\parskip}{0.5em}

% Títulos mais compactos
\titlespacing*{\section}{0pt}{0.6em}{0.3em}
\titlespacing*{\subsection}{0pt}{0.5em}{0.2em}

% Ambiente para respostas longas
\newcommand{\answerlines}[1]{% #1 = altura em cm
  \vspace*{#1}
}

% Caixa de instruções
\newenvironment{instructions}{%
  \begin{center}
    \setlength{\fboxsep}{8pt}%
    \fbox{\begin{minipage}{0.95\linewidth}\small
}{%
    \end{minipage}}%
  \end{center}
}

\begin{document}

\vspace{-1cm}

% ------------------------------------------------------------------
\section*{Parte A — Matemática como ferramenta}

\textbf{1.} Explique qual é o papel do limite no cálculo diferencial e integral. Dê um
exemplo de como ele é usado em uma situação prática. Que tipo de modelo físico pode
ser descrito usando limites?

\answerlines{3.2cm}

\textbf{2.} Explique como você entende a diferença entre \textit{derivada} e
\textit{integral}. Dê um exemplo de situação física em que cada uma aparece.

\answerlines{3.2cm}

\textbf{3.} No espaço tridimensional, o que significa dizer que $\vec{A}\cdot\vec{B}=0$?
Comente uma implicação geométrica e uma física possível.

\answerlines{3.2cm}

% ------------------------------------------------------------------
\section*{Parte B — Fundamentos de Física}

\textbf{4.} O que é um \textit{referencial} em Física? Por que precisamos escolhê-lo ao
descrever movimentos?

\answerlines{3.2cm}

\textbf{5.} Uma bola é lançada verticalmente para cima e retorna ao ponto de lançamento.
Descreva qualitativamente (sem contas) como variam, ao longo do tempo: \textit{posição},
\textit{velocidade} e \textit{aceleração}.

\answerlines{4.2cm}

\textbf{6.} Por que as \textit{leis de conservação} (por exemplo, energia, momento
linear, carga elétrica) são centrais em Física? Dê um exemplo de problema em que elas
ajudam a resolver sem muitos detalhes matemáticos.

\answerlines{3.6cm}

\textbf{7.} Diferencie uma \textit{lei física} (ex.: gravitação de Newton) de uma
\textit{definição} (ex.: velocidade média). Explique com um exemplo de cada.

\pagebreak

% ------------------------------------------------------------------
\section*{Parte C — Questões praticas}

\textbf{8. Derivada simples} \\
Calcule a derivada de
\[
  f(x) = x^2 \sin x.
\]

\vspace{3cm}

\textbf{9. Integral elementar} \\
Calcule a integral indefinida
\[
  \int (2x + 3)\, dx.
\]

\vspace{3cm}

\textbf{10. Identidade trigonométrica} \\
Mostre que (use a definição de seno e cosseno em termos dos catetos e hipotenusa para um
triângulo retângulo)
\[
  \sin^2\theta + \cos^2\theta = 1
\]
e, a partir disso, escreva $\tan^2\theta$ em função de $\sec^2\theta$.

\vspace{3cm}

\textbf{11. Manipulação algébrica/física simples} \\
Um corpo em queda livre tem posição dada por
\[
  y(t) = y_0 + v_0 t - \tfrac{1}{2}gt^2.
\]
Encontre:
\begin{enumerate}
  \item[a)] a velocidade $v(t)$;
  \item[b)] o instante em que a velocidade se anula.
\end{enumerate}




% ------------------------------------------------------------------
\section*{Parte D — Autoavaliação}

\textbf{12.} Em quais tópicos de \textit{matemática} ou \textit{física} você se sente
mais seguro? Diga por quê.

\answerlines{2.8cm}

\textbf{13.} Em quais tópicos você sente mais dificuldade ou insegurança? O que ajudaria
você a progredir?

\answerlines{2.8cm}

\textbf{14.} O que você espera aprender neste curso? Há algum tema específico de
interesse?

\answerlines{2.8cm}


\end{document}
